% Copyright 2026 MarcosHCK
%
\documentclass[12pt,a4paper]{article}
\usepackage{amsfonts}
\usepackage{amsmath}
\usepackage{amssymb}
\usepackage{bookmark}
\usepackage{fontawesome5}
\usepackage{hyperref}
\usepackage{indentfirst}
\usepackage{inputenc}[utf8]
\usepackage{minted}
\usepackage{polyglossia}
\usepackage{setspace}
\usepackage{xcolor}[dvipsnames]

\setmainlanguage{spanish}
\setotherlanguage{english}
\setstretch{1.4}

\definecolor{DarkAquamarine}{rgb}{0.18, 0.55, 0.34}
\hypersetup {
  colorlinks = true,
  linkcolor  = violet,
  urlcolor   = DarkAquamarine,
}

\begin{document}

\begin{titlepage}
  \centering
  {\huge\bfseries Reporte (semana 7+)\par}
  \vspace{1em}
  {\Large Autor: \itshape Marcos Antonio Pérez Lorenzo\par}
\end{titlepage}

\pagebreak
% Copyright 2026 MarcosHCK
%
\section{Arquitecturas generales}

\subsection{Arduino UNO}
\label{sec:genarc:arduino}

Todos los programas escritos para ejecutarse en un \textbf{Arduino UNO} R3 siguen las siguientes características:

\begin{itemize}
  \item Lenguaje: \textbf{C++}.
  \item Dependencias: compilador de C/C++ para la arquitectura AVR (específicamente el ATmega328P)
  \item Flujo: obtención de datos, cálculos, salida del programa.
\end{itemize}

\subsection{C++}
\label{sec:genarc:cplusplus}

Todos los programas escritos en \textbf{C++} siguen las siguientes características

\begin{itemize}
  \item Lenguaje: \textbf{C++}.
  \item Dependencias: librería estándar.
  \item Flujo: obtención de datos, cálculos, salida del programa.
\end{itemize}

\subsection{Python}
\label{sec:genarc:python}

Todos los programas escritos en \textbf{Python} siguen las siguientes características

\begin{itemize}
  \item Lenguaje: \textbf{Python}.
  \item Dependencias: librería estándar.
  \item Flujo: obtención de datos, cálculos, salida del programa.
\end{itemize}
% Copyright 2026 MarcosHCK
%
\section{Instrucciones de uso}

\subsection{Arduino UNO}
\label{src:genericinstrc:arduino}

Para ejecutar un sketch de Arduino, compilarlo para la placa correspondiente,
cargar el firmware y seguir la interfaz de entrada y salida que describe su
sección específica.

\subsection{C++}
\label{src:genericinstrc:cplusplus}

Todos los programas escritos en \textbf{C++} se pueden ejecutar mediante un script de \textbf{Python} intermediario para su fácil uso, por lo tanto se ejecutan de la misma forma usado para programas escritos en \textbf{Python} (descritos en \ref{src:genericinstrc:python})

\subsection{Python}
\label{src:genericinstrc:python}

Todos los programas escritos en \textbf{Python} se pueden ejecutar desde una terminal con acceso al interprete. Por ejemplo, para ejecutar el programa "clase567": en una terminal cuyo directorio de trabajo es el directorio raíz del repositorio, ejecutar "\mintinline{bash}{python3 clase567}".

\section{Programas}

Los programas y pruebas descritos en este reporte pertenecen al subproyecto
\texttt{clase01a}, escrito en Python. El código está disponible en el
\href{https://github.com/MarcosHCK/master-programming/tree/master/clase01a}
{repositorio \faGithubSquare}. El optimizador es genérico respecto al tipo de
individuo, pero su función de aptitud debe devolver un único valor real; las
estrategias de selección y la preservación de élites se configuran por separado.

Las cinco entradas siguientes documentan las estrategias y extensiones
solicitadas. Se distingue explícitamente lo que existe en el código revisado de
lo que se plantea como extensión: de ese modo, la explicación de optimización
multiobjetivo no se confunde con una capacidad ya implementada.

\subsection{Programa 1: Optimizador genético con cuatro esquemas de apareamiento}

\subsubsection{Objetivos}

El programa ejecutable \texttt{clase01a} busca maximizar la aptitud de una
población de individuos reales. El argumento \mintinline{bash}{--selection-type}
permite escoger cuatro esquemas básicos de apareamiento:

\begin{enumerate}
  \item \textbf{Mejor con todos} (\texttt{best-with-all}): el individuo mejor
    clasificado se aparea con cada uno de los demás. Si la población ordenada es
    \(r_1,r_2,\ldots,r_n\), los pares son
    \((r_1,r_2),(r_1,r_3),\ldots,(r_1,r_n)\). Este esquema genera más hijos que
    los otros; la población siguiente se limita a su tamaño configurado
    conservando los candidatos de mayor aptitud.
  \item \textbf{Aleatorio} (\texttt{random}): mezcla la población y aparea a los
    vecinos de la lista mezclada. Si el tamaño es impar, el último individuo no
    apareado pasa a la siguiente generación.
  \item \textbf{Mejor con mejor} (\texttt{elitist} en el código): ordena por
    aptitud descendente y aparea posiciones adyacentes: mejor con segundo mejor,
    tercero con cuarto, etc. El nombre \texttt{elitist} designa aquí el
    apareamiento por rango; no significa que los padres se copien sin cambios.
  \item \textbf{Mejor con peor} (\texttt{best-with-worse}): aparea extremos
    opuestos del orden: mejor con peor, segundo mejor con segundo peor, y así
    sucesivamente. Si hay un individuo central sin pareja, se conserva sin
    cruzamiento.
\end{enumerate}

\subsubsection{Tecnicismos}

El optimizador maximiza puntuaciones: una puntuación mayor siempre se considera
mejor. Para cada pareja se aplica cruzamiento binario simulado (SBX) y después
mutación gaussiana, con una probabilidad de \(0.15\) por coordenada y desviación
típica \(2.0\). El generador inicializa cada coordenada en el intervalo
\([-100,100]\); el SBX no impone límites a los hijos.

En el estado del código revisado, el individuo de ejemplo tiene dos
coordenadas \((x,y)\), pero su aptitud sigue siendo escalar:

\[
  f(x,y)=-x^2-y^2.
\]

Esta función tiene un único máximo global en \((0,0)\).

\subsubsection{Ejecución}

Desde la raíz del repositorio, por ejemplo:

\begin{minted}{bash}
python -m clase01a --selection-type best-with-all
python -m clase01a --selection-type random
python -m clase01a --selection-type elitist
python -m clase01a --selection-type best-with-worse
\end{minted}

La semilla permite repetir las decisiones aleatorias de una corrida. El
programa imprime el mejor individuo encontrado, su aptitud y las generaciones
completadas.

\subsection{Programa 2: Optimizador genético multiobjetivo y superficie de Pareto}

\subsubsection{Objetivos y fundamento}

Este apartado describe la extensión multiobjetivo solicitada y la diferencia
conceptual respecto al optimizador escalar. En un problema con varios objetivos
se evalúa cada individuo mediante un vector, por ejemplo
\(\mathbf{F}(x)=(f_1(x),f_2(x))\), en vez de una sola aptitud. Si ambos objetivos
se maximizan, un individuo \(a\) domina a \(b\) cuando no es peor en ninguno y
es estrictamente mejor en al menos uno:

\[
  f_i(a)\geq f_i(b)\quad\text{para todo }i,
  \qquad f_j(a)>f_j(b)\quad\text{para algún }j.
\]

Cuando los objetivos entran en conflicto, puede no existir un individuo que
maximice todos simultáneamente. La frontera de Pareto reúne las soluciones no
dominadas: mejorar un objetivo desde uno de esos puntos requiere empeorar al
menos otro. Por eso, una corrida ordinaria que compara individuos mediante una
sola puntuación selecciona un ganador, no la frontera completa. Si se reduce el
vector a una suma ponderada,

\[
  g_{w}(x)=w f_1(x)+(1-w)f_2(x),\qquad 0\leq w\leq1,
\]

el algoritmo maximiza \(g_w\) y devuelve una solución favorecida por ese peso.
Cambiar \(w\), la población inicial o la semilla puede conducir a otro punto
eficiente. Si varios candidatos empatan en el valor escalar, el algoritmo
conserva uno de los mejores candidatos encontrados; no devuelve por ello todos
los empates ni toda la superficie de Pareto.

\subsection{Programa 3: Algoritmo genético con elitismo añadido}

\subsubsection{Objetivos}

El parámetro \mintinline{bash}{--elite-size} conserva sin modificar los mejores
\(k\) individuos de cada generación y los copia a la siguiente. El valor
predeterminado es cero, por lo que la preservación explícita está desactivada
por omisión.

\subsubsection{Diferencia respecto al apareamiento elitista}

El elitismo añadido es independiente de la estrategia \texttt{elitist} del
programa 1. Esa estrategia únicamente aparea individuos consecutivos del
ranking; sus hijos pueden reemplazar a los padres. En cambio,
\texttt{elite\_size} reserva puestos para individuos de la población actual
copiados literalmente, independientemente del método de apareamiento elegido.
Por ejemplo, se puede combinar \texttt{best-with-worse} con dos élites:

\begin{minted}{bash}
python -m clase01a --selection-type best-with-worse --elite-size 2
\end{minted}

El algoritmo también lleva un registro del mejor individuo encontrado a lo
largo de toda la corrida. Ese registro histórico no es lo mismo que preservar
ese individuo en cada población; la garantía de conservación generacional la
proporciona \texttt{elite\_size}.

\subsection{Programa 4: Selección por ruleta y por torneo}

\subsubsection{Selección por ruleta}

La opción \texttt{roulette} elige cada padre mediante muestreo con probabilidad
proporcional a un peso derivado de su aptitud. Si las puntuaciones son
negativas, primero se desplazan para obtener pesos no negativos; si todas son
iguales, se usan pesos uniformes. El muestreo es con reemplazo, por lo que un
mismo individuo puede ser seleccionado más de una vez en una generación. El
método ofrece una preferencia probabilística por candidatos mejores, sin
restringir la selección exclusivamente al primero del ranking.

\subsubsection{Selección por torneo}

La opción \texttt{tournament} toma tres candidatos al azar para cada elección
de padre y escoge el de mayor aptitud. Los dos padres de una pareja se
seleccionan por separado, por lo que un individuo puede volver a participar.
Para poblaciones menores que tres, el torneo usa todos los candidatos
disponibles. Aumentar el tamaño del torneo normalmente refuerza la presión de
selección; en esta implementación el tamaño está fijado en tres.

\subsubsection{Ejecución}

Ambas opciones usan los mismos argumentos de población, generaciones, semilla
y tamaño de élite que las demás estrategias:

\begin{minted}{bash}
python -m clase01a --selection-type roulette --elite-size 0
python -m clase01a --selection-type tournament --elite-size 2
\end{minted}

Las seis opciones disponibles en el CLI son \texttt{elitist}, \texttt{random},
\texttt{best-with-all}, \texttt{best-with-worse}, \texttt{tournament} y
\texttt{roulette}.

\subsection{Programa 5: Animación de las estrategias de optimización}

\subsubsection{Objetivos y funcionamiento}

El subcomando \texttt{animate} registra poblaciones durante la optimización y
genera una gráfica por generación capturada. Incluye la generación inicial,
cada \(N\)-ésima generación (\texttt{--generation-divider}) y siempre la
generación final. Las gráficas se ensamblan como GIF con Pillow o como WebM con
\texttt{ffmpeg}; los cuadros PNG pueden guardarse en un directorio o eliminarse
al terminar. La representación muestra la curva de aptitud, los individuos de
la población y el mejor de la generación.

El script \path{clase01a/animate.sh} recorre las seis estrategias dos veces:
primero con \texttt{--elite-size 0} y después con \texttt{--elite-size 2}.
Configura 200 generaciones, un cuadro cada cinco generaciones, 100 ms por
cuadro y límites horizontales \([-120,120]\). Así se preparan doce salidas en
\path{clase01a/animations/}:

\begin{itemize}
  \item \texttt{elitist.webm}, \texttt{elitist-elite.webm};
  \item \texttt{random.webm}, \texttt{random-elite.webm};
  \item \texttt{best-with-all.webm}, \texttt{best-with-all-elite.webm};
  \item \texttt{best-with-worse.webm}, \texttt{best-with-worse-elite.webm};
  \item \texttt{tournament.webm}, \texttt{tournament-elite.webm};
  \item \texttt{roulette.webm}, \texttt{roulette-elite.webm}.
\end{itemize}

Se puede generar una animación individual directamente, por ejemplo:

\begin{minted}{bash}
python -m clase01a animate --selection-type roulette --elite-size 2 \
  --generations 200 --generation-divider 5 --frame-duration 100 \
  --output clase01a/animations/roulette-elite.webm --x-limits -120 120
\end{minted}

\end{document}